\documentclass[12pt]{article}

\usepackage{graphicx}
\usepackage[hypertexnames=false,colorlinks=true,breaklinks]{hyperref}
\usepackage{tikz}
\usepackage{verbatim}
\usepackage{natbib}
\usepackage{apalike}
\usepackage{amsmath,amssymb}
\usepackage{endfloat}

\newcommand{\estNumber}{82227365}

\title{Detection of replay on Emmett's data with Linear Dynamical Systems}
\author{Joaqu\'{i}n Rapela\thanks{j.rapela@ucl.ac.uk}}

\begin{document}

\maketitle

\section{Introduction}

Here we use a Linear Dynamical System
model (LDS)\footnote{\url{https://github.com/joacorapela/ssm}} to detect replay
events in of Emmett's switching-task recordings~\citep{thompsonEtAl24}.

\section{Methods}

\subsection{Data}

We analyzed the spiking activity of all striatal neurons (neuron indices 0-207)
in electrophysiological recording
\texttt{EJT178\_implant1/recording6\_29\-03\-2022}.
%
Figure~\ref{fig:spikeTimes} plots a sample of spike times used in this
analysis.

\begin{figure}
    \centering
    \href{http://www.gatsby.ucl.ac.uk/~rapela/emmett/reports/lds_20260307/figures/spike_times_start4045.00_end4345.00.html}{\includegraphics[width=6in]{../../../figures/EJT178_implant1/recording6_29-03-2022/spike_times_start4045.00_end4345.00.png}}

    \caption{Example spikes times used in this analysis. We used recording
    session \texttt{EJT178\_implant1/recording6\_29-03-2022} and neurons
    indices 0 to 207 in the striatum.
    %
    Click on the image to get its interactive version.
    %
    }

    \label{fig:spikeTimes}
\end{figure}

\subsection{LDS model}

The LDS model estimates $K$ latent variables, $x_0[n],\ldots,x_{K-1}[n]$,
$n\in\{0,\ldots,N-1\}$.  We
concatenated the latent variables into a latent vector $\mathbf{x}[n]$:

\begin{align}
    \mathbf{x}[n]=\left[\begin{array}{c}
        x_0[n]\\
                            \vdots\\
                            x_{K-1}[n]
                        \end{array}\right]
    \label{eq:latentsVector}
\end{align}

The latents evolve in time following a first-order Markov process with additive
Gaussian noise:

\begin{align*}
    \mathbf{x}[n]=A\mathbf{x}[n-1]+\mathbf{v}[n]\quad\text{with}\ \mathbf{v}[n]\sim\mathcal{N}(0, Q)
\end{align*}

\noindent starting from
$\mathbf{x}_0\sim\mathcal{N}(\mathbf{x}_0|\mathbf{m}_0,V_0)$.

The latent variables are combined linearly to approximate the firing rate of
the recorded neurons (bin size 200~ms), $\mathbf{y}[n]$, with a constant offset,
$\mathbf{a}$, and an additive noise, $\mathbf{w}[n]$.

\begin{align*}
    \mathbf{y}[n]=\mathbf{a}+C\mathbf{x}[n]+\mathbf{w}[n]\quad\text{with}\ \mathbf{w}[n]\sim\mathcal{N}(0, R)
\end{align*}

We estimated an LDS models with $K=5$ latent variables
(Eq.~\ref{eq:latentsVector}; the selection of the number of latent variables
was arbitrary, and we should perform model selection in later iterations).

We estimate  the model parameter $\{\mathbf{m}_0, V_0, \mathbf{u}, A, Q,
\mathbf{a}, Z, R\}$ using the expectation maximisation algorithm
(\href{https://github.com/joacorapela/lds_neuralLatents_striatum/blob/master/code/scripts/doEstimateEM.py}{doEstimateEM.py}).
%
We then used the Kalman Filter algorithm, with these parameters, to infer the
latents variables, $\mathbf{x}[n]$ in Equation~\ref{eq:latentsVector}
(\href{https://github.com/joacorapela/lds_neuralLatents_striatum/blob/master/code/scripts/doKalmanFilteringWithEstimatedParams.py}{doKalmanFilteringWithEstimatedParams.py}).
%
% We also forecasted firing rates of individual neurons
% (\href{https://github.com/joacorapela/lds_neuralLatents_striatum/blob/master/code/scripts/doPlotObservationsForecasts.py}{doPlotObservationsForecasts.py}).

\section{Results}

\begin{figure} \centering
    \href{http://www.gatsby.ucl.ac.uk/~rapela/emmett/reports/lds_20260307/figures/42904130_state_filtered_from4045.00_to4345.00.html}{\includegraphics[width=6in]{../../../figures/EJT178_implant1/recording6_29-03-2022/42904130_state_filtered_from4045.00_to4345.00.png}}

    \caption{Orthonormalised latent variables ($\mathbf{x}(t)$ in
    Eq.~\ref{eq:latentsVector}) inferred from the data in the behavioral period
    used to estimate the models parameters (between 4045 and 4345 sec).
    %
    The colored vertical lines indicate `poke\_in` times, with green, red,
    blue, yellow, purple, cyan and magenta corresponding to ports 1 to 7,
    respectively.
    %
    The correct sequence is $2\rightarrow 1\rightarrow 6\rightarrow
    3\rightarrow 7$ corresponding to
    $\text{red}\rightarrow\text{green}\rightarrow\text{blue}\rightarrow\text{cyan}\rightarrow\text{magenta}$.
    %
    Click on the image to get its interactive version.
    %
    }

    \label{fig:inferredLatentsTrain}
\end{figure}

\begin{figure}
    \centering

    \href{http://www.gatsby.ucl.ac.uk/~rapela/emmett/reports/lds_20260307/figures/42904130_hmm_from4045.00_to4345.00_mean.html}{\includegraphics[width=6in]{../../../figures/EJT178_implant1/recording6_29-03-2022/42904130_hmm_from4045.00_to4345.00_mean.png}}

    \caption{HMM emission parameters. Inferred means and standard deviations (error
    bars) of the Gaussian emission model for each latent (i.e., observation)
    and for each HMM state.}

    \label{fig:hmm_emision_params}
\end{figure}

\begin{figure}
    \centering

    \href{http://www.gatsby.ucl.ac.uk/~rapela/emmett/reports/lds_20260307/figures/42904130_hmm_from4045.00_to4345.00_transition.html}{\includegraphics[width=6in]{../../../figures/EJT178_implant1/recording6_29-03-2022/42904130_hmm_from4045.00_to4345.00_transition.png}}

    \caption{Inferred HMM transition matrix.}

    \label{fig:hmm_transition_matrix}
\end{figure}

\begin{figure}
    \centering

    \href{http://www.gatsby.ucl.ac.uk/~rapela/emmett/reports/lds_20260307/figures/train42904130_test42904130_true_vs_inferred_hmm_most_prob_state_seq_trainFrom4045.00_trainTo4345.00_testFrom4045.00_testTo4345.00.html}{\includegraphics[width=6in]{../../../figures/EJT178_implant1/recording6_29-03-2022/train42904130_test42904130_true_vs_inferred_hmm_most_prob_state_seq_trainFrom4045.00_trainTo4345.00_testFrom4045.00_testTo4345.00.png}}
    \caption{True (blue) versus estimated (red) behavioral states, when
    estimation and inference was done on the same data segment of the
    behavioral period. The LDS and
    HMM models were estimated with data between 4045 to 4345~sec. The LDS model
    inferred latents and the HMM inferred behavioral states from the same
    data segment.}

    \label{fig:trueVsInferredStatesTest4045_4345}
\end{figure}

\begin{figure}
    \centering
    \href{http://www.gatsby.ucl.ac.uk/~rapela/emmett/reports/lds_20260307/figures/12157029_state_filtered_from4345.00_to4645.00.html}{\includegraphics[width=6in]{../../../figures/EJT178_implant1/recording6_29-03-2022/12157029_state_filtered_from4345.00_to4645.00.png}}

    \caption{Orthonormalised latent variables inferred from neural data in the
    behavioral period not used to train LDS or HMM model parameters (in the
    behavioral segment between 4343 and 4645~sec).}

    \label{fig:inferredLatentsTestDuringBehavior}
\end{figure}

\begin{figure}
    \centering

    \href{http://www.gatsby.ucl.ac.uk/~rapela/emmett/reports/lds_20260307/figures/train42904130_test12157029_true_vs_inferred_hmm_most_prob_state_seq_trainFrom4045.00_trainTo4345.00_testFrom4345.00_testTo4645.00.html}{\includegraphics[width=6in]{../../../figures/EJT178_implant1/recording6_29-03-2022/train42904130_test12157029_true_vs_inferred_hmm_most_prob_state_seq_trainFrom4045.00_trainTo4345.00_testFrom4345.00_testTo4645.00.png}}

    \caption{True (blue) versus estimated (red) behavioral states, when
    estimation and inference were done on different data segments of the
    behavioral period. The LDS and
    HMM models were estimated with data between 4045 to 4345~sec. The LDS model
    inferred latents and the HMM inferred behavioral states from data between
    4345 to 4645~sec.}

    \label{fig:trueVsInferredStatesTest4345_4645}
\end{figure}

\begin{figure}
    \centering

    \href{http://www.gatsby.ucl.ac.uk/~rapela/emmett/reports/lds_20260307/figures/21792914_state_filtered_from6000.00_to7454.00.html}{\includegraphics[width=6in]{../../../figures/EJT178_implant1/recording6_29-03-2022/21792914_state_filtered_from6000.00_to7454.00.png}}
    \caption{Orthonormalised latent variables inferred from data in a
    non-behavioral block immediately following behavior between 6000 and
    7454~sec.}

    \label{fig:inferredLatentsTestDuringPostBehavior}
\end{figure}

\begin{figure}
    \centering

    \href{http://www.gatsby.ucl.ac.uk/~rapela/emmett/reports/lds_20260307/figures/train42904130_test21792914_inferred_hmm_most_prob_state_seq_trainFrom4045.00_trainTo4345.00_testFrom6000.00_testTo7454.00.html}{\includegraphics[width=6in]{../../../figures/EJT178_implant1/recording6_29-03-2022/train42904130_test21792914_inferred_hmm_most_prob_state_seq_trainFrom4045.00_trainTo4345.00_testFrom6000.00_testTo7454.00.png}}

    \caption{True (blue) versus estimated (red) behavioral states, when
    estimation was done with data in the behavioral period and inference was
    done with data in a non-behavioral period immediately following behavior.
    The LDS and HMM models were estimated with data between 4045 to 4345~sec.
    The LDS model inferred latents and the HMM inferred behavioral states from
    data between 6000 to 7454~sec.}

    \label{fig:trueVsInferredStatesTest6000_7454}
\end{figure}

\bibliographystyle{apalike}
\bibliography{replay.bib}

\end{document}
